\documentclass[11pt]{article}
\usepackage{amsmath,amssymb,amsthm}

\newtheorem{theorem}{Theorem}[section]
\newtheorem{lemma}[theorem]{Lemma}
\newtheorem{proposition}[theorem]{Proposition}
\theoremstyle{remark}
\newtheorem{remark}[theorem]{Remark}

\newcommand{\R}{\mathbb{R}}
\newcommand{\Rpos}{\mathbb{R}_{>0}}

\title{A short note on three elementary inequalities}
\author{Benchmark manuscript for the TeX proof-check kit}
\date{}

\begin{document}
\maketitle

\begin{abstract}
We record three elementary inequalities for real numbers that are used
repeatedly in later work.  The note is self-contained except for one
classical inequality, which we quote from the literature.
\end{abstract}

\section{Notation}

Throughout, $\R$ denotes the real numbers and $\Rpos$ the positive real
numbers.  For real numbers $x,y,z$ we write
\[
  S(x,y,z) = \frac{x}{y+z} + \frac{y}{z+x} + \frac{z}{x+y}
\]
whenever the three denominators are nonzero.

\section{Results}

\begin{lemma}\label{lem:amgm}
For all real numbers $a$ and $b$ we have
\[
  ab \le \frac{a^2+b^2}{2},
\]
with equality if and only if $a = b$.
\end{lemma}

\begin{proof}
Since $(a-b)^2 \ge 0$, expanding the square gives $a^2 - 2ab + b^2 \ge 0$,
that is, $2ab \le a^2 + b^2$.  Dividing by $2$ gives the inequality.
Moreover, equality holds if and only if $(a-b)^2 = 0$, that is, if and only
if $a = b$.
\end{proof}

\begin{proposition}\label{prop:nesbitt}
For all $x,y,z \in \Rpos$ we have
\[
  S(x,y,z) \ge \frac{3\,(xy+yz+zx)}{2\,(x^2+y^2+z^2)}.
\]
\end{proposition}

\begin{proof}
By Nesbitt's inequality \cite[Theorem~1]{Nesbitt1903}, $S(x,y,z) \ge 3/2$
for all $x,y,z \in \Rpos$.  Applying Lemma~\ref{lem:amgm} to the pairs
$(x,y)$, $(y,z)$ and $(z,x)$ and adding the three inequalities gives
\[
  xy + yz + zx \le x^2 + y^2 + z^2 .
\]
Since $x^2+y^2+z^2 > 0$, dividing by it yields
$\frac{xy+yz+zx}{x^2+y^2+z^2} \le 1$, and hence
$\frac{3\,(xy+yz+zx)}{2\,(x^2+y^2+z^2)} \le \frac{3}{2} \le S(x,y,z)$.
\end{proof}

\begin{lemma}\label{lem:quot}
For every real number $x > -1$ we have
\[
  \frac{x^2}{1+x} \ge x - 1 .
\]
\end{lemma}

\begin{proof}
Since $1+x > 0$, multiplying both sides by $1+x$ preserves the inequality, so the claim
is equivalent to $x^2 \ge (x-1)(x+1) = x^2 - 1$, which is obvious.
\end{proof}

\begin{remark}
Lemma~\ref{lem:quot} is used in the companion paper only for $x \ge 0$.
\end{remark}

\begin{thebibliography}{9}
\bibitem{Nesbitt1903}
A.~M. Nesbitt, Problem 15114, \emph{Educational Times} \textbf{55} (1903), 37.
\end{thebibliography}

\end{document}
